% === ЗАДАЧА 1 ===
% Тег: #ОПТИКА #ОТРАЖЕНИЕ #КИМ_13
\begin{taskbasic}[code={\label{task:geom_opt_mirror_angle_34}}]
Луч света падает на плоское зеркало. Угол между падающим и отражённым лучами равен $34^\circ$. Чему равен угол между падающим лучом и зеркалом? Ответ выразите в градусах ($^\circ$).

\kim{13}
\end{taskbasic}
\answer{73}

% === ЗАДАЧА 2 ===
% Тег: #ОПТИКА #ОТРАЖЕНИЕ #КИМ_13
\begin{taskbasic}[code={\label{task:geom_opt_mirror_angle_54}}]
Луч света падает на плоское зеркало. Угол между падающим и отражённым лучами равен $54^\circ$. Чему равен угол между отражённым лучом и зеркалом? Ответ выразите в градусах ($^\circ$).

\kim{13}
\end{taskbasic}
\answer{63}

% === ЗАДАЧА 3 ===
% Тег: #ОПТИКА #ОТРАЖЕНИЕ #КИМ_13
\begin{taskbasic}[code={\label{task:geom_opt_mirror_angle_15_between}}]
Луч света падает на плоское зеркало. Угол падения равен $15^\circ$. Чему равен угол $\gamma$ между падающим и отражённым лучами? Ответ выразите в градусах ($^\circ$).

\kim{13}
\end{taskbasic}
\answer{30}

% === ЗАДАЧА 4 ===
% Тег: #ОПТИКА #ОТРАЖЕНИЕ #КИМ_13
\begin{taskbasic}[code={\label{task:geom_opt_mirror_angle_70_between}}]
Луч света падает на плоское зеркало. Угол между отражённым лучом и зеркалом равен $70^\circ$. Чему равен угол $\gamma$ между падающим и отражённым лучами? Ответ выразите в градусах ($^\circ$).

\kim{13}
\end{taskbasic}
\answer{40}

% === ЗАДАЧА 5 ===
% Тег: #ОПТИКА #ОТРАЖЕНИЕ #КИМ_13
\begin{taskbasic}[code={\label{task:geom_opt_mirror_refl_15_surface}}]
Луч света падает на плоское зеркало. Угол отражения равен $15^\circ$. Определите угол между падающим лучом и плоскостью зеркала. Ответ выразите в градусах ($^\circ$).

\kim{13}
\end{taskbasic}
\answer{75}

% === ЗАДАЧА 6 ===
% Тег: #ОПТИКА #ОТРАЖЕНИЕ #КИМ_13
\begin{taskbasic}[code={\label{task:geom_opt_mirror_refl_30_between}}]
Луч света падает на плоское зеркало. Угол отражения равен $30^\circ$. Определите угол между падающим и отражённым лучами. Ответ выразите в градусах ($^\circ$).

\kim{13}
\end{taskbasic}
\answer{60}

% === ЗАДАЧА 7 ===
% Тег: #ОПТИКА #ОТРАЖЕНИЕ #КИМ_13
\begin{taskbasic}[code={\label{task:geom_opt_mirror_refl_25_surface}}]
Луч света падает на плоское зеркало. Угол отражения равен $25^\circ$. Определите угол между падающим лучом и поверхностью зеркала. Ответ выразите в градусах ($^\circ$).

\kim{13}
\end{taskbasic}
\answer{65}

% === ЗАДАЧА 8 ===
% Тег: #ОПТИКА #ОТРАЖЕНИЕ #КИМ_13
\begin{taskbasic}[code={\label{task:geom_opt_mirror_inc_20_between}}]
Луч света падает на плоское зеркало. Угол падения равен $20^\circ$. Определите угол между падающим и отражённым лучами. Ответ выразите в градусах ($^\circ$).

\kim{13}
\end{taskbasic}
\answer{40}

% === ЗАДАЧА 9 ===
% Тег: #ОПТИКА #ПРЕЛОМЛЕНИЕ #ВОЛНЫ #КИМ_15
\begin{taskbasic}[code={\label{task:geom_opt_wave_air_plexiglas}}]
Световой пучок переходит из воздуха в плексиглас (см. рисунок). Что происходит при этом с длиной световой волны и скоростью её распространения?

Для каждой величины определите соответствующий характер изменения:
1) увеличивается
2) уменьшается
3) не изменяется

\nopagebreak\smallskip
\begin{center}
\begin{tikzpicture}[scale=1.0, every node/.style={font=\footnotesize}]
    \draw[xstep=1, ystep=1, gray!30, thin] (0,0) grid (4,3);
    
    % Граница раздела (воздух / плексиглас)
    \fill[gray!30] (0,0) rectangle (4,1.5);
    \draw[thick, black] (0,1.5) -- (4,1.5);
    
    % Нормаль
    \draw[dashed, black, thick] (2,0) -- (2,3);
    
    % Падающий луч
    \draw[-{Stealth[scale=0.8]}, thick, black] (0.8,2.7) -- (1.4,2.1);
    \draw[thick, black] (1.4,2.1) -- (2.0,1.5);
    \draw (2.0,2.1) arc (90:124:0.6);
    \node[above left] at (1.8,2.0) {$\alpha$};
    
    % Преломлённый луч
    \draw[-{Stealth[scale=0.8]}, thick, black] (2.0,1.5) -- (2.35,0.75);
    \draw[thick, black] (2.35,0.75) -- (2.7,0);
    \draw (2.0,1.0) arc (-90:-63:0.5);
    \node[below right] at (1.9,1.1) {$\beta$};
    
    \node[below left] at (0,0) {0};
\end{tikzpicture}
\end{center}

\kim{15}
\end{taskbasic}
\answer{22}

% === ЗАДАЧА 10 ===
% Тег: #ОПТИКА #ПРЕЛОМЛЕНИЕ #ВОЛНЫ #КИМ_15
\begin{taskbasic}[code={\label{task:geom_opt_wave_air_glycerin}}]
Плоская световая волна переходит из воздуха в глицерин (см. рисунок). Что происходит при этом переходе с частотой электромагнитных колебаний в световой волне и скоростью их распространения?

Для каждой величины определите соответствующий характер изменения:
1) увеличивается
2) уменьшается
3) не изменяется

\nopagebreak\smallskip
\begin{center}
\begin{tikzpicture}[scale=1.0, every node/.style={font=\footnotesize}]
    \draw[xstep=1, ystep=1, gray!30, thin] (0,0) grid (4,3);
    
    % Граница раздела
    \fill[gray!30] (0,0) rectangle (4,1.5);
    \draw[thick, black] (0,1.5) -- (4,1.5);
    
    % Нормаль
    \draw[dashed, black, thick] (2,0) -- (2,3);
    
    % Падающий луч
    \draw[-{Stealth[scale=0.8]}, thick, black] (0.8,2.7) -- (1.4,2.1);
    \draw[thick, black] (1.4,2.1) -- (2.0,1.5);
    \draw (2.0,2.1) arc (90:124:0.6);
    \node[above left] at (1.8,2.0) {$\alpha$};
    
    % Преломлённый луч
    \draw[-{Stealth[scale=0.8]}, thick, black] (2.0,1.5) -- (2.35,0.75);
    \draw[thick, black] (2.35,0.75) -- (2.7,0);
    \draw (2.0,1.0) arc (-90:-63:0.5);
    \node[below right] at (1.9,1.1) {$\beta$};
    
    \node[below left] at (0,0) {0};
\end{tikzpicture}
\end{center}

\kim{15}
\end{taskbasic}
\answer{32}

% === ЗАДАЧА 11 ===
% Тег: #ОПТИКА #ОТРАЖЕНИЕ #ЗЕРКАЛО_ПОВОРОТ #КИМ_13
\begin{taskbasic}[code={\label{task:geom_opt_mirror_rotation_gamma_between}}]
Угол падения луча света на горизонтальное плоское зеркало равен $25^\circ$. Каким будет угол $\gamma$, образованный падающим и отражённым лучами, если повернуть зеркало на $10^\circ$ так, как показано на рисунке? Ответ выразите в градусах ($^\circ$).

\nopagebreak\smallskip
\begin{center}
\begin{tikzpicture}[scale=1.0, every node/.style={font=\footnotesize}]
    \draw[xstep=1, ystep=1, gray!30, thin] (0,0) grid (5,3);
    
    % Поверхность горизонтальная и под углом 10 град
    \fill[gray!30] (1.0,0.8) -- (4.0,0.8) -- (4.0,0.6) -- (1.0,0.6) -- cycle;
    \draw[thick, black] (1.0,0.8) -- (4.0,0.8);
    \draw[dashed, black, thick] (1.0,0.8) -- (4.0,0.2);
    \draw (3.5,0.8) arc (0:-11:1.0);
    \node[right] at (3.6,0.6) {$10^\circ$};
    
    % Нормаль и лучи
    \draw[dashed, gray, thick] (2.5,0.8) -- (2.5,2.8);
    \draw[-{Stealth[scale=0.8]}, thick, black] (1.0,2.8) -- (2.5,0.8);
    \draw[-{Stealth[scale=0.8]}, thick, black] (2.5,0.8) -- (4.2,2.8);
    
    \draw (2.5,1.8) arc (90:126:1.0);
    \node[left] at (2.4,2.0) {$25^\circ$};
    
    \draw (2.0,1.5) arc (126:48:0.8);
    \node[above] at (2.8,1.4) {$\gamma$};
    
    \node[below left] at (0,0) {0};
\end{tikzpicture}
\end{center}

\kim{13}
\end{taskbasic}
\answer{70}

% === ЗАДАЧА 12 ===
% Тег: #ОПТИКА #ОТРАЖЕНИЕ #ЗЕРКАЛО_ПОВОРОТ #КИМ_13
\begin{taskbasic}[code={\label{task:geom_opt_mirror_rotation_gamma_refl}}]
Угол падения луча света на горизонтальное плоское зеркало равен $25^\circ$. Каким будет угол отражения $\gamma$, если повернуть зеркало на $15^\circ$ так, как показано на рисунке? Ответ выразите в градусах ($^\circ$).

\nopagebreak\smallskip
\begin{center}
\begin{tikzpicture}[scale=1.0, every node/.style={font=\footnotesize}]
    \draw[xstep=1, ystep=1, gray!30, thin] (0,0) grid (5,3);
    
    % Поверхность горизонтальная и под углом 15 град
    \fill[gray!30] (1.0,0.8) -- (4.0,0.8) -- (4.0,0.6) -- (1.0,0.6) -- cycle;
    \draw[thick, black] (1.0,0.8) -- (4.0,0.8);
    \draw[dashed, black, thick] (1.0,0.8) -- (4.0,0.1);
    \draw (3.3,0.8) arc (0:-15:0.8);
    \node[right] at (3.5,0.5) {$15^\circ$};
    
    % Нормаль и лучи
    \draw[dashed, gray, thick] (2.5,0.8) -- (2.5,2.8);
    \draw[-{Stealth[scale=0.8]}, thick, black] (1.0,2.8) -- (2.5,0.8);
    \draw[-{Stealth[scale=0.8]}, thick, black] (2.5,0.8) -- (4.0,2.8);
    
    \draw (2.5,1.8) arc (90:126:1.0);
    \node[left] at (2.4,2.0) {$25^\circ$};
    
    \draw (2.5,1.4) arc (90:53:0.6);
    \node[right] at (2.5,1.6) {$\gamma$};
    
    \node[below left] at (0,0) {0};
\end{tikzpicture}
\end{center}

\kim{13}
\end{taskbasic}
\answer{40}

\begin{center}
\begin{tikzpicture}[scale=1.0, transform shape, >=Stealth]

  % Параметры
  \def\R{3.5}       % Радиус лимба (шкалы)
  \def\Rg{2.0}      % Радиус стеклянного полуцилиндра

  % --- Внешний круг шкалы ---
  \draw[line width=1.0pt] (0,0) circle (\R);

  % --- Деления шкалы (каждые 1°, 5° и 10°) ---
  \foreach \a in {0,1,...,359} {
    % Исключаем главные деления (0, 90, 180, 270), так как они отрисовываются отдельно
    \ifnum\a=0\else\ifnum\a=90\else\ifnum\a=180\else\ifnum\a=270\else
      \pgfmathparse{mod(\a,10)==0 ? 1 : (mod(\a,5)==0 ? 2 : 3)}
      \ifnum\pgfmathresult=1
        \draw[thick] (\a:\R) -- (\a:{\R-0.3});
      \else
        \ifnum\pgfmathresult=2
          \draw[thin] (\a:\R) -- (\a:{\R-0.2});
        \else
          \draw[very thin, gray!80] (\a:\R) -- (\a:{\R-0.12});
        \fi
      \fi
    \fi\fi\fi\fi
  }

  % --- Удлинённые ключевые метки для 0° и 90° (все четыре направления) ---
  \foreach \a in {0, 90, 180, 270} {
    \draw[line width=1.2pt] (\a:\R) -- (\a:{\R-0.5});
  }

  % Обозначения 90
  \node[font=\large\bfseries] at (90:2.7) {90};
  \node[font=\large\bfseries, rotate=180] at (270:2.7) {90};

  % --- Стеклянный полуцилиндр ---
  \draw[line width=1.0pt, fill=white] (-\Rg, 0) -- (\Rg, 0) arc (0:-180:\Rg) -- cycle;

  % Метки 0 (повёрнуты на 90° и сдвинуты к внешним делениям)
  \node[font=\large\bfseries, rotate=90] at (180:2.5) {0};
  \node[font=\large\bfseries, rotate=90] at (0:2.5) {0};

  % --- Лучи света ---
  
  % 1. Падающий луч
  \draw[line width=1.1pt, ->] (160:\R) -- (0,0);

  % Двойная стрелка-указатель у входа падающего луча
  \begin{scope}[shift={(160:{\R+0.1})}, rotate=-20]
    \draw[line width=0.9pt, fill=white] (-0.3, 0.15) -- (0, 0) -- (-0.3, -0.15) -- (-0.18, 0) -- cycle;
  \end{scope}

  % 2. Отраженный луч
  \draw[line width=1.1pt, ->] (0,0) -- (20:\R);

  % 3. Преломленный луч
  \draw[line width=1.1pt, ->] (0,0) -- (310:\R);

\end{tikzpicture}
\end{center}

\begin{center}
\begin{tikzpicture}[scale=1.0, transform shape, >=Stealth]

  % Параметры прямоугольного треугольника ABC
  \def\a{5.0}  % Длина катета AC (основание)
  \def\b{2.8}  % Длина катета BC (высота)

  % Координаты вершин
  \coordinate (A) at (0, 0);
  \coordinate (C) at (\a, 0);
  \coordinate (B) at (\a, \b);

  % --- Гипотенуза AB (жирная линия) ---
  \draw[line width=2.2pt] (A) -- (B);

  % --- Катеты AC и BC (обычные линии) ---
  \draw[thick] (A) -- (C) -- (B);

  % --- Прямой угол C ---
  \draw[thick] (\a-0.3, 0) -- (\a-0.3, 0.3) -- (\a, 0.3);

  % --- Дуга и обозначение угла alpha ---
  \draw[thick] (0.8, 0) arc (0:29.25:0.8);
  \node[font=\small\bfseries\itshape] at (0.55, 0.16) {$\alpha$};

  % --- Подписи вершин ---
  \node[left=2pt, font=\large\bfseries\itshape] at (A) {$A$};
  \node[above right=1pt, font=\large\bfseries\itshape] at (B) {$B$};
  \node[below right=1pt, font=\large\bfseries\itshape] at (C) {$C$};

  % --- Вертикальный луч света (стрелка снизу вверх) ---
  \def\xRay{3.5} % Координата x входящего луча

  % Стрелка падающего луча
  \draw[line width=1.3pt, ->] (\xRay, -1.0) -- (\xRay, 0);

  % Прямой угол падения на грань AC
  \draw[thick] (\xRay, -0.3) -- (\xRay+0.3, -0.3) -- (\xRay+0.3, 0);

\end{tikzpicture}
\end{center}